\documentclass[10pt,a4paper]{article}

%double si on veut une version prof et une version élève. simple si on veut une seule version.
\newcommand{\buildmode}{simple}

\InputIfFileExists{version.tex}{}{}
% Version (définie par le script)
\providecommand{\version}{eleve}

\usepackage{feuilleinterro}

%%%à modifier à chaque fois

\renewcommand{\niveau}{1M}
\renewcommand{\chapitre}{Proportions, pourcentages}
\renewcommand{\numchap}{0.2}
\renewcommand{\theme}{Statistiques}

%%%titre "CORRECTION" au lieu de "INTERROGATION", sans les champs Nom/Prénom
\renewcommand{\interroversion}[1]{
    \noindent\textbf{\niveau}

    \vspace{-0.5cm}
    \begin{center}
        \textbf{\large CORRECTION -- \chapitre}
    \end{center}

    \vspace{0.1cm}
}

\begin{document}

% =========================================================
% PAGE 1 : VERSION A
% =========================================================

\interroversion{A}

\textbf{Exercice 1 -- Calcul avec des fractions \hfill /3}

\begin{enumerate}
\item \(A = \dfrac{6}{7}\times\dfrac{8}{9}\)

\[
\begin{aligned}
A &= \frac{6\times8}{7\times9}\\
  &= \frac{48}{63}\\
  &= \frac{16}{21}.
\end{aligned}
\]

\item \(B = \dfrac{5}{6}-\dfrac{1}{3}\)

\[
\begin{aligned}
B &= \frac{5}{6}-\frac{2}{6}\\
  &= \frac{3}{6}=\frac{1}{2}.
\end{aligned}
\]
\end{enumerate}

\textbf{Exercice 2 -- Proportion, pourcentage et écriture décimale \hfill /3}

Dans un lycée de 200 élèves, 50 sont internes.

La proportion d'élèves internes est \(\dfrac{50}{200}=\dfrac{1}{4}\).

Sous forme de pourcentage : \(\dfrac{1}{4}=\dfrac{25}{100}=\textcolor{blue}{25\%}\).

Sous forme de nombre décimal : \(\textcolor{blue}{0{,}25}\).

\textbf{Exercice 3 -- Tableau croisé \hfill /4}

\begin{enumerate}
\item Compléter les deux cases manquantes du tableau. \hfill /2

Total des filles (somme sur la ligne) : \(50+30=\textcolor{blue}{80}\).

Effectif des garçons externes (soustraction sur la ligne) : \(120-70=\textcolor{blue}{50}\).

\begin{center}
\renewcommand{\arraystretch}{1.5}
\begin{tabular}{|l|c|c|c|}
\hline
 & \textbf{Interne} & \textbf{Externe} & \textbf{Total} \\
\hline
\textbf{Filles} & 50 & 30 & \textcolor{blue}{80} \\
\hline
\textbf{Garçons} & 70 & \textcolor{blue}{50} & 120 \\
\hline
\textbf{Total} & 120 & 80 & 200 \\
\hline
\end{tabular}
\end{center}

\item Quelle est la proportion de filles parmi l'ensemble des élèves du lycée ? \hfill /1

\(\dfrac{80}{200}=\dfrac{2}{5}=\textcolor{blue}{0{,}4=40\%}\).

\item Parmi les filles, quelle est la proportion d'élèves internes ? \hfill /1

\(\dfrac{50}{80}=\dfrac{5}{8}=\textcolor{blue}{0{,}625=62{,}5\%}\).
\end{enumerate}

\newpage

% =========================================================
% PAGE 2 : VERSION B
% =========================================================

\interroversion{B}

\textbf{Exercice 1 -- Calcul avec des fractions \hfill /3}

\begin{enumerate}
\item \(A = \dfrac{5}{9}\times\dfrac{7}{10}\)

\[
\begin{aligned}
A &= \frac{5\times7}{9\times10}\\
  &= \frac{35}{90}\\
  &= \frac{7}{18}.
\end{aligned}
\]

\item \(B = \dfrac{7}{8}-\dfrac{1}{4}\)

\[
\begin{aligned}
B &= \frac{7}{8}-\frac{2}{8}\\
  &= \frac{5}{8}.
\end{aligned}
\]
\end{enumerate}

\textbf{Exercice 2 -- Proportion, pourcentage et écriture décimale \hfill /3}

Un club de sport compte 200 adhérents, dont 50 pratiquent la natation.

La proportion d'adhérents pratiquant la natation est \(\dfrac{50}{200}=\dfrac{1}{4}\).

Sous forme de pourcentage : \(\dfrac{1}{4}=\dfrac{25}{100}=\textcolor{blue}{25\%}\).

Sous forme de nombre décimal : \(\textcolor{blue}{0{,}25}\).

\textbf{Exercice 3 -- Tableau croisé \hfill /4}

\begin{enumerate}
\item Compléter les deux cases manquantes du tableau. \hfill /2

Total des femmes (somme sur la ligne) : \(60+40=\textcolor{blue}{100}\).

Effectif des hommes en abonnement mensuel (soustraction sur la ligne) : \(150-90=\textcolor{blue}{60}\).

\begin{center}
\renewcommand{\arraystretch}{1.5}
\begin{tabular}{|l|c|c|c|}
\hline
 & \textbf{Annuel} & \textbf{Mensuel} & \textbf{Total} \\
\hline
\textbf{Femmes} & 60 & 40 & \textcolor{blue}{100} \\
\hline
\textbf{Hommes} & 90 & \textcolor{blue}{60} & 150 \\
\hline
\textbf{Total} & 150 & 100 & 250 \\
\hline
\end{tabular}
\end{center}

\item Quelle est la proportion de femmes parmi l'ensemble des adhérents ? \hfill /1

\(\dfrac{100}{250}=\dfrac{2}{5}=\textcolor{blue}{0{,}4=40\%}\).

\item Parmi les femmes, quelle est la proportion d'abonnements annuels ? \hfill /1

\(\dfrac{60}{100}=\dfrac{3}{5}=\textcolor{blue}{0{,}6=60\%}\).
\end{enumerate}

\newpage

% =========================================================
% PAGE 3 : VERSION C
% =========================================================

\interroversion{C}

\textbf{Exercice 1 -- Calcul avec des fractions \hfill /3}

\begin{enumerate}
\item \(A = \dfrac{8}{11}\times\dfrac{5}{6}\)

\[
\begin{aligned}
A &= \frac{8\times5}{11\times6}\\
  &= \frac{40}{66}\\
  &= \frac{20}{33}.
\end{aligned}
\]

\item \(B = \dfrac{7}{9}-\dfrac{1}{3}\)

\[
\begin{aligned}
B &= \frac{7}{9}-\frac{3}{9}\\
  &= \frac{4}{9}.
\end{aligned}
\]
\end{enumerate}

\textbf{Exercice 2 -- Proportion, pourcentage et écriture décimale \hfill /3}

Une entreprise emploie 200 salariés, dont 50 travaillent à temps partiel.

La proportion de salariés à temps partiel est \(\dfrac{50}{200}=\dfrac{1}{4}\).

Sous forme de pourcentage : \(\dfrac{1}{4}=\dfrac{25}{100}=\textcolor{blue}{25\%}\).

Sous forme de nombre décimal : \(\textcolor{blue}{0{,}25}\).

\textbf{Exercice 3 -- Tableau croisé \hfill /4}

\begin{enumerate}
\item Compléter les deux cases manquantes du tableau. \hfill /2

Total des femmes (somme sur la ligne) : \(50+50=\textcolor{blue}{100}\).

Effectif des hommes non-cadres (soustraction sur la ligne) : \(150-70=\textcolor{blue}{80}\).

\begin{center}
\renewcommand{\arraystretch}{1.5}
\begin{tabular}{|l|c|c|c|}
\hline
 & \textbf{Cadre} & \textbf{Non-cadre} & \textbf{Total} \\
\hline
\textbf{Femmes} & 50 & 50 & \textcolor{blue}{100} \\
\hline
\textbf{Hommes} & 70 & \textcolor{blue}{80} & 150 \\
\hline
\textbf{Total} & 120 & 130 & 250 \\
\hline
\end{tabular}
\end{center}

\item Quelle est la proportion de femmes parmi l'ensemble des salariés ? \hfill /1

\(\dfrac{100}{250}=\dfrac{2}{5}=\textcolor{blue}{0{,}4=40\%}\).

\item Parmi les femmes, quelle est la proportion de cadres ? \hfill /1

\(\dfrac{50}{100}=\dfrac{1}{2}=\textcolor{blue}{0{,}5=50\%}\).
\end{enumerate}

\newpage

% =========================================================
% PAGE 4 : VERSION D
% =========================================================

\interroversion{D}

\textbf{Exercice 1 -- Calcul avec des fractions \hfill /3}

\begin{enumerate}
\item \(A = \dfrac{9}{10}\times\dfrac{6}{7}\)

\[
\begin{aligned}
A &= \frac{9\times6}{10\times7}\\
  &= \frac{54}{70}\\
  &= \frac{27}{35}.
\end{aligned}
\]

\item \(B = \dfrac{5}{8}-\dfrac{1}{4}\)

\[
\begin{aligned}
B &= \frac{5}{8}-\frac{2}{8}\\
  &= \frac{3}{8}.
\end{aligned}
\]
\end{enumerate}

\textbf{Exercice 2 -- Proportion, pourcentage et écriture décimale \hfill /3}

Un festival a accueilli 200 festivaliers, dont 50 avaient moins de 18 ans.

La proportion de festivaliers de moins de 18 ans est \(\dfrac{50}{200}=\dfrac{1}{4}\).

Sous forme de pourcentage : \(\dfrac{1}{4}=\dfrac{25}{100}=\textcolor{blue}{25\%}\).

Sous forme de nombre décimal : \(\textcolor{blue}{0{,}25}\).

\textbf{Exercice 3 -- Tableau croisé \hfill /4}

\begin{enumerate}
\item Compléter les deux cases manquantes du tableau. \hfill /2

Total des Pass 3 jours (somme sur la ligne) : \(30+70=\textcolor{blue}{100}\).

Effectif des Pass 1 jour âgés de 18 ans et plus (soustraction sur la ligne) : \(150-50=\textcolor{blue}{100}\).

\begin{center}
\renewcommand{\arraystretch}{1.5}
\begin{tabular}{|l|c|c|c|}
\hline
 & \textbf{Moins de 18 ans} & \textbf{18 ans et plus} & \textbf{Total} \\
\hline
\textbf{Pass 3 jours} & 30 & 70 & \textcolor{blue}{100} \\
\hline
\textbf{Pass 1 jour} & 50 & \textcolor{blue}{100} & 150 \\
\hline
\textbf{Total} & 80 & 170 & 250 \\
\hline
\end{tabular}
\end{center}

\item Quelle est la proportion de festivaliers ayant un pass 3 jours parmi l'ensemble des festivaliers ? \hfill /1

\(\dfrac{100}{250}=\dfrac{2}{5}=\textcolor{blue}{0{,}4=40\%}\).

\item Parmi les festivaliers ayant un pass 3 jours, quelle est la proportion de moins de 18 ans ? \hfill /1

\(\dfrac{30}{100}=\dfrac{3}{10}=\textcolor{blue}{0{,}3=30\%}\).
\end{enumerate}

\end{document}
